% ===========================================================================
\section{Study Area and Data}
\label{sec:data}
% ===========================================================================

\subsection{Population, which is the activity datum}

Every category in this sector is driven by population. Waste generation is a per-capita rate
applied to urban and rural residents; the organic load entering the wastewater system is a
per-capita BOD applied to the same population split by sanitation pathway. There is no
tonnage measurement and no metered flow anywhere in the chain, so the population panel is the
activity data in a way that is unusual for an emissions inventory.

The panel is built at district level and interpolated between censuses, covering
\num{134} districts, with the 2023 census as its reference---the census's full frame
including the Islamabad Capital Territory, Gilgit-Baltistan and Azad Kashmir, a
deliberately wider base than the \num{116} four-province districts the agriculture
sector allocates to, whose activity statistics are provincial in scope. It resolves to
\SI{97}{} million urban and \SI{150}{} million rural residents, \SI{247}{} million in total.
The urban/rural split is load-bearing twice over: urban residents generate more waste per head
and are far more likely to have their waste collected, and the two effects run in opposite
directions for emissions, since collected waste goes to a decaying landfill while uncollected
waste burns.

Building the panel at district level rather than nationally is what makes the provincial
allocation in Section~\ref{sec:results} a real geography rather than a national figure split
by a population share, and it is the same panel that drives the sector's air-pollutant layer.

\subsection{Waste generation and collection}

Per-capita generation is \SI{0.453}{\kilo\gram} per person per day urban and
\SI{0.283}{\kilo\gram} rural, from the Pak-EPA eight-city survey \citep{pakepa1996}, the
rates the \citet{kaza2018} Pakistan row re-bins. Across the panel that is
\SI{31.56}{\mega\tonne} of municipal solid waste a year.

Collection is expressed as its complement, the uncollected fraction: \SI{30}{\percent} urban
and \SI{70}{\percent} rural, which resolves to \SI{50.3}{\percent} of national generation
collected. Of the uncollected remainder, \SI{60}{\percent} is openly burned; the rest is
dumped, tipped into watercourses or otherwise disposed of without combustion and emits nothing
this category reports.

These three parameters---two generation rates, two uncollected fractions and one burned
fraction---are the softest inputs in the sector, and Section~\ref{sec:limitations} says so.
They are also the only place where a judgement moves emissions \emph{between} categories
rather than changing the total, which is a property of the partition rather than of the
parameters.

\subsection{Waste composition}

Open burning needs the carbon content of what burns. Dry matter is \SI{70}{\percent} of wet
weight; total carbon is \SI{0.2195}{\tonne} per wet tonne and the fossil fraction of it is
\SI{0.0757}{\tonne} per wet tonne, so roughly a third of the carbon in Pakistani municipal
waste is fossil in origin. That third is what produces reportable \COtwo{}; the biogenic
remainder is a memo item, following IPCC Volume~5 Chapter~5 \S5.1.

\subsection{Sanitation}

The wastewater category needs the population split across sanitation pathways, because the
methane correction factor is a property of the pathway rather than of the person. Shares come
from the WHO/UNICEF Joint Monitoring Programme \citep{jmp2023}, each year taking its own:
sewered, septic, latrine and open defecation,
each carrying its own correction factor (\num{0.5}, \num{0.5}, \num{0.1} and \num{0}
respectively).

Organic load is \SI{40}{} g BOD per person per day, the regional default covering Pakistan,
with an industrial-and-commercial correction of \num{1.25} for load discharged to the domestic
system from small commercial premises. Nitrous oxide is computed from protein consumption
rather than from BOD, which is why it is invariant under both of the sensitivities reported
in Section~\ref{sec:results}.

\subsection{What this inventory has not obtained, and what it would change}

Nothing compiled here reports landfill tonnages at facility level, methane recovery at any
site, industrial solid-waste quantities, or industrial wastewater organic loading. Three of
those four have leads that have not been followed: the Urban Unit has mapped solid-waste
services and facilities across Punjab's tehsils, PBS's district-level living-standards survey
records household disposal route---municipal collection, private, public bins, roadside,
waterway, open space---and PCRWR has published organic loading for named Pakistani drains.
None would give complete national coverage, and each would replace a national assumption with
district or facility evidence for the population it covers. The first two
would replace modelled disposal with measured disposal; the second two are whole categories
this inventory reports as not estimated. Section~\ref{sec:limitations} takes each in turn.
